
\chapter*{前言}
\addcontentsline{toc}{chapter}{前言}
\markboth{前言}{前言}

\begingroup\small

自从 \BookRecordIndex{T0003}$J$-全纯曲线理论于 1985 年 Gromov 的论文中诞生以来，它一直对辛拓扑学家具有重大意义。它的应用包括辛拓扑中的许多关键结果，也是\BookRecordIndex{T0132}弗洛尔同调创立的主要启发之一。它还引起了数学物理学家的关注，因为它为定义\BookRecordIndex{T0143}格罗莫夫--威滕不变量和\BookRecordIndex{T0229}量子上同调提供了自然的背景，而这两者构成镜像对称猜想中所谓的 A 侧。物理学的见解反过来又激发了许多引人入胜的进展，例如揭示了可积系统理论与\BookRecordIndex{T0143}格罗莫夫--威滕不变量之间迄今仍少有人理解的联系。

几年前，本书的两位作者写过一部介绍这一领域的著作，阐释 \BookRecordIndex{T0003}$J$-全纯曲线理论的主要技术步骤。本书起初是作为那本书的第二版构思的，但计划很快扩展了。这个领域发展如此迅速，以至于几乎没有时间巩固其基础；而基础中又涉及许多分析上的微妙之处，这已成为相当大的障碍。因此，本书的首要目标是完整而严格、细致地确立本学科的基本定理。我们也希望本书能作为辛拓扑当前研究工作的入门读物。这两个目标当然在一定程度上相互冲突，因此本书不同部分各有侧重。

我们选择着重建立理论的基础，而不试图详尽论述近来的众多进展。因此，我们将范围限定于亏格为零的曲线（不过，除球面外也讨论圆盘）。一个更重大的限制是只处理\BookRecordIndex{T0058}半正情形；在这一情形中，可以借助\BookRecordIndex{T0337}伪循环定义\BookRecordIndex{T0143}格罗莫夫--威滕不变量。这样做的主要原因是，一般情形的最优框架（其中需要构造\BookRecordIndex{T0370}虚模循环）尚未形成。我们不愿拼凑一个足以应付许多应用、却不能适用于\BookRecordIndex{T0352}辛场论等更广泛语境的定义，而是决定展示一种更简单、更几何的方法所能取得的成果。另一方面，我们对基本\BookRecordIndex{T0256}粘合定理给出了十分详尽的证明。它是此后所有\BookRecordIndex{T0370}虚模循环研究的分析基础，也是证明量子乘法结合律的关键环节。书中还有五个篇幅较长的附录，主题从\BookRecordIndex{T0380}隐函数定理、\BookRecordIndex{T0329}椭圆正则性和 \BookRecordIndex{T0221}黎曼--罗赫定理等标准结果，到亏格为零的\BookRecordIndex{T0339}稳定曲线\BookRecordIndex{T0246}模空间的结构、四维 \BookRecordIndex{T0003}$J$-全纯曲线的\BookRecordIndex{T0186}交点正性等较少为人熟知的内容。对于应用，我们也采取同样的方法：完整证明基础性结果，并通过若干关键例子和丰富的参考文献说明较近的进展。

本书的编排使各主题按逻辑顺序展开。第 2 至 5 章建立 $J$-全纯球面和圆盘所需的基础弗雷德霍姆理论与\BookRecordIndex{T0192}紧性结果；第 6 章引入在\BookRecordIndex{T0058}半正流形上定义\BookRecordIndex{T0145}格罗莫夫--威滕伪循环所需的概念；第 7 章是定义这些不变量的枢纽；其后各章讨论各种应用。由于第 2 至 6 章的细节多到初读时不可能全部消化，我们写了导论性的第 1 章，概述理论轮廓并介绍主要定义。它也充当本书的详细导读，指出关键论证所在，以及理解各种例子所需的背景细节可在何处找到。每一章也有介绍主要内容的导言，以帮助已有一定知识的读者把握方向。我们尽可能使各节、各章彼此独立；因此，读者遇到看来过于技术性的部分，不妨跳过。

我们希望第 1 章（辅以第 2 至 6 章适当部分）能为初学者提供足够的基础，使他们理解第 7 章的主要定义。以下简述其余各章的内容。在第 7.1 节定义基本不变量后（第 7.2 节和第 7.3 节还有重要的补充思想），第 7.4 节讨论射影空间中\BookRecordIndex{T0386}有理曲线这一基本例子。本章最后讨论亏格为零的\BookRecordIndex{T0143}格罗莫夫--威滕不变量的 Kontsevich--Manin 公理，并由此推导 Kontsevich 的优美递推公式，用于计算射影平面中次数为 $d$ 的\BookRecordIndex{T0386}有理曲线数目。

第 8 章建立以黎曼曲面为底的\BookRecordIndex{T0200}局部哈密顿纤维化理论，并说明如何计数这类\BookRecordIndex{T0279}纤维化的截面。这使我们能够定义任意亏格的\BookRecordIndex{T0143}格罗莫夫--威滕不变量（但定义域的复结构须固定）。它也为一些重要应用提供背景，例如 Gromov 关于\BookRecordIndex{T0355}辛非球面流形上的每个哈密顿系统都有一条 1-周期轨道的结果（见定理 9.1.1），以及关于哈密顿辛同胚群的结果：第 9.6 节初步介绍 Hofer 几何，第 11.4 节和第 12.5 节则讨论 \BookRecordIndex{T0047}Seidel 表示。

第 9 章介绍 \BookRecordIndex{T0003}$J$-全纯曲线技术在辛几何中的若干主要应用。除上述例子和对\BookRecordIndex{T0219}拉格朗日子流形基本性质的讨论外，本章还完整证明 McDuff 关于有理曲面和直纹曲面这两类辛四维流形结构的结果，以及 Gromov 关于射影平面和两个二维球面之积的辛同胚群的结果。

另一项主要应用——\BookRecordIndex{T0229}量子上同调——还需要一项深刻的分析技术，即\BookRecordIndex{T0254}粘合。最早的严格\BookRecordIndex{T0254}粘合论证分别归功于 Floer（在略为容易的\BookRecordIndex{T0132}弗洛尔同调情形中）和 Ruan--Tian（在与\BookRecordIndex{T0230}量子同调有关的情形中）。第 10 章提出另一种或许更简便的\BookRecordIndex{T0254}粘合方法，并据此证明\BookRecordIndex{T0058}半正流形上的\BookRecordIndex{T0143}格罗莫夫--威滕不变量的\BookRecordIndex{T0125}分裂公理。

在此基础上，第 11 章定义\BookRecordIndex{T0229}量子上同调，并阐释由此产生的若干结构，例如\BookRecordIndex{T0144}格罗莫夫--威滕势和 \BookRecordIndex{T0020}Frobenius 流形。从第 11.3 节的例子可以看出，辛拓扑正是在这里与代数几何、共形场论、镜像对称和可积系统等其他领域发生最深刻的联系。读完第 7 章后，读者应能阅读这一章。最后，第 12 章是一篇综述，勾勒弗洛尔理论的主要轮廓，略去其分析基础。它以几何方法解释弗洛尔理论与\BookRecordIndex{T0229}量子上同调的关系，并指出进一步发展的方向，包括分析方面的涡旋方程和几何方面的 \BookRecordIndex{T0014}Donaldson 量子范畴。

本书有五个附录。前三个附录建立经典线性椭圆算子理论的基础，而第 3、4 章将推广这一理论：附录 A 讨论弗雷德霍姆理论与\BookRecordIndex{T0053}巴拿赫流形上的\BookRecordIndex{T0380}隐函数定理，附录 B 讨论 \BookRecordIndex{T0319}索伯列夫空间与\BookRecordIndex{T0329}椭圆正则性，附录 C 讨论有边界黎曼曲面的 \BookRecordIndex{T0221}黎曼--罗赫定理。附录 D 为第 5 章提供背景；它不用通常的代数几何方法，而用\BookRecordIndex{T0185}交比说明亏格为零的\BookRecordIndex{T0339}稳定曲线的 Grothendieck--Knudsen \BookRecordIndex{T0246}模空间结构。附录 E 由我们与 Laurent Lazzarini 合作撰写，其中完整证明四维流形中 \BookRecordIndex{T0003}$J$-全纯曲线的\BookRecordIndex{T0186}交点正性和\BookRecordIndex{T0054}伴随不等式。Lazzarini 提供了这一附录的初稿及完整证明，此后我们共同完善了文字。附录 E 的结果为有理曲面和直纹曲面这两类辛四维流形的结构定理奠定基础。

希望用本书作为研究生课程基础的教师，必须明确决定自己想讲授何种课程。根据我们的经验，一个学期不可能既证明全部主要分析技术，又介绍有趣的例子。第 1 章更详细地说明了一种可能的安排：集中学习第 1 章（辅以第 2 章的少量内容）、第 3 章至第 3.3 节（或许再补充附录 B、C 的一些分析知识）、第 4.2 节关于最小能量球面的基本\BookRecordIndex{T0192}紧性结果，以及第 6 章中精挑细选的部分（\BookRecordIndex{T0337}伪循环的定义），然后进入第 7.1 节。接下来，可以直接学习第 9 章的一些几何应用（例如证明\BookRecordIndex{T0120}非挤压定理，或第 9.4 节关于辛四维流形的部分结果）；也可以讨论第 7.5 节中\BookRecordIndex{T0143}格罗莫夫--威滕不变量的 Kontsevich--Manin 公理，再转入第 11 章建立\BookRecordIndex{T0229}量子上同调。这样安排的用意是：先熟悉主要的分析框架，证明一些基本技术，再通过讨论一组应用，把这些技术放入更广阔的背景中。

上述安排或许仍过于雄心勃勃，但也有办法缩短预备部分。例如，学完第 2 至 4 章的基础材料（以及第 8 章的相关部分）后，可以直接讨论第 9 章的许多应用，而不必涉及第 5、6、7 章。因为只考虑这些应用中最简单的情形、而不证明其最一般的形式时，相关\BookRecordIndex{T0246}模空间是紧的，所以无须正式定义\BookRecordIndex{T0143}格罗莫夫--威滕不变量，也能理解这些结果。另一种选择是，若课程以\BookRecordIndex{T0229}量子上同调为目标，可以不证明弗雷德霍姆理论的结果，而着重解释第 4、5 章中的一些\BookRecordIndex{T0192}紧性（\BookRecordIndex{T0243}冒泡）结果。这些内容与\BookRecordIndex{T0340}稳定映射\BookRecordIndex{T0246}模空间的研究能够很好地结合，进而自然引出 Kontsevich--Manin 公理。

如上所述，除非讲课速度快得惊人，或几乎不涉及应用，一门初级课程不可能既讲弗雷德霍姆理论，又深入掌握\BookRecordIndex{T0192}紧性证明的分析细节，更不用说详述\BookRecordIndex{T0254}粘合的全部细节。尽管\BookRecordIndex{T0256}粘合定理的证明大体上需要与第 3 章相同的分析背景，证明\BookRecordIndex{T0257}粘合映射的满射性仍有赖于第 4 章最深刻的结果（小能量长圆柱的行为），并依赖若干技术性估计。我们撰写\BookRecordIndex{T0254}粘合一章，力求使读者能够把握构造的轮廓和主要分析思想（这些内容在第 10.5 节作了概述）。因此，对于分析基础较深的读者，也可以用第 3、4、10 章组织一门课程，并以第 9 或第 11 章的一些例子作为学习动机。

尽管本书篇幅很长，主题却如此丰富，以至于不可能面面俱到。我们在全书各处给出了许多参考文献。下面列举几部相关领域的著作；读者既可直接参阅，也可利用其中进一步的文献线索：Cox--Katz \srcite{76} 论述镜像对称与代数几何，Donaldson \srcite{87} 论述\BookRecordIndex{T0132}弗洛尔同调与规范理论，Manin \srcite{286} 论述 \BookRecordIndex{T0020}Frobenius 流形与\BookRecordIndex{T0229}量子上同调，Polterovich \srcite{330} 论述辛同胚群的几何，而 Eliashberg--Givental--Hofer \srcite{101} 的论文论述\BookRecordIndex{T0352}辛场论。

本书历经多年才得以完成；若无许多同事慷慨分享见解与知识，它不可能问世。尤其是 Coates、Givental、Lalonde、Lazzarini、Polterovich、Popescu、Robbin、Ruan 和 Seidel，都为本书的不同部分提供了关键帮助。我们还要感谢指出各种排印错误和不准确之处的众多学生与其他读者，尤其感谢 Eduardo Gonzalez、Sam Lisi、Jake Solomon 和 Fabian Ziltener 对细节一丝不苟的关注。

\begin{flushright}
Dusa McDuff 与 Dietmar Salamon，2003 年 12 月
\end{flushright}
\endgroup
